\section{Method}
\label{sec:method}

\subsection{Corpora}

DAIGT V2~\cite{Thedrcat2023daigtv2} contains 44,868 argumentative student essays, 27,371 human-written and 17,497 machine-generated by a mixture of 2023-era systems including GPT-3.5, Llama-2, Mistral-7B and Falcon-180B, assembled for a public detection competition~\cite{Kaggle2023llmdetect}. The human side derives from the PERSUADE corpus of student argumentative writing~\cite{Crossley2023persuade}. HC3~\cite{Guo2023close} contains question-answer pairs from five English sources, contrasting human answers with a single generator, GPT-3.5-Turbo-0301, over 85,449 rows~\cite{HelloSimpleAI2023hc3}. The two differ in almost every respect that matters here, which is why they are a useful pair. DAIGT V2 has many generators, one genre and long documents. HC3 has one generator, five source domains and short ones.

Both corpora were class-balanced by downsampling the majority class, giving 34,994 DAIGT V2 rows and 53,806 HC3 rows. Balancing is not cosmetic. It fixes the score of a degenerate all-one-class predictor at 0.333 weighted F1 rather than at something respectable, which is what makes the label-free control in Section~\ref{sec:labelfree} readable. Training partitions are balanced to within 0.002 of parity.

\begin{table}[!ht]
\centering
\caption{Corpus statistics after balancing. Duplication is measured over the unbalanced corpus by exact match on whitespace-normalised, lowercased text. Leakage is the share of test documents whose exact text also appears in training or validation, measured on the balanced sample under each split rule. The group-aware rule is the one every number in this paper uses.}
\label{tab:corpora}
\footnotesize
\setlength{\tabcolsep}{3.5pt}
\begin{tabular}{@{}lrr@{}}
\toprule
& DAIGT V2 & HC3 \\
\midrule
Rows, as published & 44,868 & 85,449 \\
Rows, balanced & 34,994 & 53,806 \\
Train / validation / test & 25,196 & 38,785 \\
& 2,800 & 4,289 \\
& 6,998 & 10,732 \\
\midrule
Duplicate rows & 6 (0.01\%) & 6,118 (7.16\%) \\
Duplicate groups & 6 & 5,986 \\
Cross-label duplicate texts & 0 & 0 \\
Collection artefact rows & none found & 483 \\
\midrule
Test leakage, group-aware split & 0 (0.00\%) & 0 (0.00\%) \\
Test leakage, naive stratified & 0 (0.00\%) & 570 (5.30\%) \\
\bottomrule
\end{tabular}
\end{table}

A duplicate-detection pass over the unbalanced corpora finds DAIGT V2 essentially free of internal duplication, six rows in six groups, with no text appearing under both labels. HC3 carries 6,118 duplicate rows in 5,986 groups, plus 476 rows labelled machine-generated and 7 labelled human that are collection artefacts, typically stored API failure messages. Table~\ref{tab:corpora} gives the figures.

\subsection{Partitioning}

That duplication is why the split is group-aware, and the choice is worth more than it sounds. Documents are grouped by an MD5 hash of their whitespace-normalised, lowercased text, and whole groups are assigned to a single partition, so no two documents with identical normalised content are placed on opposite sides of a boundary. The partition is 72/8/20 rather than 80/10/10, because a fifth is taken for test first and a tenth of the remaining four-fifths for validation.

Measured directly, the group-aware rule leaks no HC3 test document into training or validation, 0 of 10,732, whereas a plain stratified split of the same balanced sample leaks 570 of 10,762, which is 5.30\%. DAIGT V2 leaks nothing either way. No number in this paper is inflated by memorised duplicates. The split is built once and reused unchanged by every model and every training seed, so a comparison between two models is never also a comparison between two evaluation sets.

\subsection{Models and configurations}

Five model families are evaluated. Three are classical, Naive Bayes, logistic regression and a linear support vector machine, each run under bag-of-words and TF-IDF for six classical configurations per dataset. Earlier work at full scale reported only the winning representation per model, so we report the complete grid. Two are transformers, \texttt{bert-base-uncased}~\cite{Devlin2019bert} and \texttt{microsoft/deberta-v3-base}~\cite{He2023debertav3}, fine-tuned with AdamW~\cite{Loshchilov2019adamw} over a grid of learning rate in $\{2\mathrm{e}{-5}, 3\mathrm{e}{-5}\}$, batch size in $\{16, 32\}$ and weight decay in $\{0.01, 0.1\}$, which is sixteen configurations per dataset, with the operating point selected on validation weighted F1 and early stopping at patience two.

DeBERTa-v3 is the variant chosen for three reasons that matter to this study rather than for its leaderboard position. Its disentangled attention separates content from position~\cite{He2021deberta}, which is the relevant inductive bias when the question is how much of a corpus is separable by form. Its ELECTRA-style pretraining~\cite{Clark2020electra} is more sample-efficient at a comparable parameter budget. And its SentencePiece tokeniser~\cite{Kudo2018sentencepiece} encodes leading whitespace, whereas BERT's WordPiece~\cite{Wu2016google} does not, which makes the pair a controlled contrast on exactly the cue Section~\ref{sec:tokenisation} examines. Implementations are HuggingFace Transformers~\cite{Wolf2020transformers} and scikit-learn~\cite{Pedregosa2011scikit}.

Classical models receive text through the normalisation pipeline described below, which lowercases, removes every character outside the basic Latin alphabet, drops English stopwords and lemmatises. Transformers receive raw text, whitespace-normalised only, truncated to 128 tokens. That asymmetry is deliberate, since stripping punctuation and casing removes signal a pretrained transformer can use, but it is also a confound for any classical-versus-transformer comparison, and Section~\ref{sec:budget} measures it rather than assuming it away.

\subsection{The surface-content decomposition}
\label{sec:decomp-def}

\begin{figure*}[!t]
\centering
\includegraphics[width=\textwidth]{figures/fig_pipeline.pdf}
\caption{The decomposition. One balanced, duplicate-group-aware partition feeds three arms. The surface arm reads 47 orthographic statistics and no word identity. The content arm reads an unfiltered bag-of-words after punctuation, casing and non-ASCII characters have been removed. The full arm is a fine-tuned transformer on raw text and is reported as a reference point, not as a matched arm. Surface and content share one classifier family so they are directly comparable, and every comparison between them is made with a paired test on the same documents. The length control at the bottom closes the one channel both primary arms can otherwise read. Counts on the diagram are read from the artefacts the experiment wrote.}
\label{fig:pipeline}
\end{figure*}

The central instrument of this paper separates how much of a benchmark's human-versus-machine separability is available from surface form alone and how much from content alone. Fig.~\ref{fig:pipeline} shows the design.

Write $x$ for a document and $y \in \{0, 1\}$ for its label, with $1$ denoting machine-generated. Let $\mathcal{D}_{\text{tr}}$ and $\mathcal{D}_{\text{te}}$ be the fixed training and test partitions. Two maps are defined on the raw text. The surface map $\phi_{\mathrm{S}}$ sends a document to $\mathbb{R}^{47}$, a vector of orthographic statistics described in Section~\ref{sec:features}, and reads no word identity whatsoever. The content map $\phi_{\mathrm{C}}$ first applies the normalisation
\begin{equation}
\label{eq:norm}
\nu(x) = \mathrm{collapse}\Big(\mathrm{map}\big(\mathrm{lower}(x),\, \sigma\big)\Big),
\end{equation}
where $\mathrm{lower}$ lowercases, $\mathrm{map}$ applies the character rule $\sigma$ that keeps a character when it is a lowercase Latin letter or whitespace and replaces it with a space otherwise, and $\mathrm{collapse}$ reduces runs of whitespace to a single space. The normalised text is then sent to raw bag-of-words counts over the training vocabulary. Punctuation, casing and non-ASCII characters cannot survive $\nu$, so they cannot enter the content arm. Word identity never enters $\phi_{\mathrm{S}}$.

Each arm fits one logistic regression on the same training partition,
\begin{equation}
\label{eq:arm}
h_{a}  =  \arg\min_{h} \, \sum_{(x,y) \in \mathcal{D}_{\text{tr}}} \ell\big(h(\phi_{a}(x)), y\big)  +  \tfrac{1}{2C}\lVert h \rVert_2^2 ,
\end{equation}
for $a \in \{\mathrm{S}, \mathrm{C}\}$, with $\ell$ the logistic loss and $C$ fixed at the library default for both arms. Sharing the classifier family and the regularisation strength is what makes the two arms comparable, since any difference between them is then a difference in what the representation carries rather than in how it is fitted.

The quantity we report is the error of each arm on the held-out partition,
\begin{equation}
\label{eq:err}
\varepsilon_{a}  =  \frac{1}{|\mathcal{D}_{\text{te}}|}\sum_{(x,y)\in\mathcal{D}_{\text{te}}} \mathbf{1}\big[h_{a}(\phi_{a}(x)) \neq y\big],
\end{equation}
and the separability gap between them,
\begin{equation}
\label{eq:gap}
\Delta  =  \varepsilon_{\mathrm{S}} - \varepsilon_{\mathrm{C}} .
\end{equation}
A gap near zero says orthography alone reaches the same accuracy as content alone on that corpus. A large positive gap says content carries what orthography does not. The surface arm's own error $\varepsilon_{\mathrm{S}}$ answers a second and more direct question, namely how much of the benchmark a detector could pass without reading language at all.

Two properties of $\Delta$ should be stated plainly. It is a comparison between two specific arms and not a partition of the total separability, so a richer surface featurisation would lower $\varepsilon_{\mathrm{S}}$ and a stronger content model would lower $\varepsilon_{\mathrm{C}}$. Both arms therefore bound rather than decompose in the strict sense, and $\Delta$ should be read as a corpus-level comparison under a fixed pair of instruments. It is also not a claim that either corpus is clean, since $\varepsilon_{\mathrm{S}}$ is far below chance on both.

The \emph{full} arm is the fine-tuned transformer on raw text. It differs from the other two in model family, in parameter count and in the amount of text it reads, so it is reported as a reference point and we draw no conclusion from a transformer-to-classical difference alone without first matching the text budget.

\subsection{The length channel and how it is closed}
\label{sec:lengthdef}

One channel survives into both arms as defined above. The surface arm reads character, word and sentence counts directly, and raw bag-of-words rows sum to document length, so $\phi_{\mathrm{C}}$ can recover it too. Length is a documented confound on both corpora~\cite{Baidya2026detecting}, so we close it rather than assume it away.

On the surface side we drop the five features whose magnitude scales with document size, leaving 42 rates and ratios, and write $\phi_{\mathrm{S}}^{-L}$ for the reduced map. Every dropped feature has a rate twin in the same vector, so the cue itself is not removed along with the length. On the content side we normalise each row to unit total and then rescale by the mean training document length $\bar{n}$,
\begin{equation}
\label{eq:l1}
\phi_{\mathrm{C}}^{-L}(x)  =  \bar{n} \cdot \frac{\phi_{\mathrm{C}}(x)}{\lVert \phi_{\mathrm{C}}(x) \rVert_1} .
\end{equation}
The rescaling is not cosmetic. Rows that sum to one shrink every feature by roughly $\bar{n}$, so at a fixed $C$ the effective penalty becomes far heavier and the arm collapses for reasons of scale rather than of length. We report both variants in Section~\ref{sec:length} and read only the rescaled one, and we treat the unrescaled arm as an instrument artefact rather than a result.

We also fit a \emph{length-only} arm on the three pure document-size features, so the strength of the shared channel is measured directly instead of being inferred from what removing it costs the other two.

\subsection{The surface feature set}
\label{sec:features}

\begin{table}[!ht]
\centering
\caption{The 47 surface features, by family. None reads word identity. The five features marked as size-scaling are the ones the length control removes, and each has a rate twin that remains. Two features are constant on the DAIGT V2 test split and one on HC3, which is recorded because a feature with no variance contributes nothing and should not be counted as evidence of coverage.}
\label{tab:features}
\footnotesize
\setlength{\tabcolsep}{3.5pt}
\begin{tabular}{@{}lcl@{}}
\toprule
Family & Count & Notes \\
\midrule
Space before punctuation & 2 & raw count is size-scaling \\
Document size & 3 & characters, words, sentences \\
Word shape & 1 & mean word length \\
Casing & 3 & uppercase, capitalised, all-caps \\
Non-ASCII and emoji & 2 & emoji count is size-scaling \\
Digits & 1 & rate per 1,000 characters \\
Whitespace layout & 2 & double spaces, newlines \\
Sentence structure & 1 & mean sentence length \\
Per-character punctuation rates & 32 & one per punctuation mark \\
\midrule
Total & 47 & 5 size-scaling, 42 rates and ratios \\
\bottomrule
\end{tabular}
\end{table}

Table~\ref{tab:features} lists the families. The set is deliberately plain. Every feature is a count or a rate computable in one pass over the raw string, with no model, no lexicon and no language identification, which is what makes the measurement cheap enough to run on any new corpus before committing to it.

\begin{figure*}[!t]
\centering
\includegraphics[width=\textwidth]{figures/fig_surface_correlation.pdf}
\caption{Spearman correlation among the 47 surface features on each test split, with rows and columns ordered by average-linkage clustering on $1-|\rho|$ so correlated blocks are adjacent. Features the length control removes are outlined and labelled in green. The set is close to unredundant on both corpora, with mean absolute off-diagonal correlation 0.062 on DAIGT V2 and 0.084 on HC3, and only 10 of 1,081 pairs on DAIGT V2 and 14 of 1,081 on HC3 exceeding 0.5 in absolute value. The visible blocks are the ones a reader would predict, namely the bracket and parenthesis rates moving together and the document-size features moving with the sentence count. Features with no variance on a test split have undefined correlation and are drawn as zero.}
\label{fig:corr}
\end{figure*}

Fig.~\ref{fig:corr} shows how much these features overlap. They are close to unredundant, which matters for two reasons. It means the surface arm's performance is not one cue counted 47 times, and it means the length control removes a small and identifiable part of the vector rather than most of it in disguise.

\subsection{Metrics}

We report weighted F1 as the headline metric, because it is what the prior work on these corpora reports, alongside error rate because F1 differences compress at this accuracy ceiling and a change from 0.9916 to 0.9972 reads as small while the error rate it corresponds to falls by a factor of three. For class $c$ with precision $P_c$ and recall $R_c$, support $n_c$ and $N=\sum_c n_c$,
\begin{equation}
\label{eq:f1}
\mathrm{F1}_c = \frac{2 P_c R_c}{P_c + R_c}, \quad
\mathrm{F1}_{\mathrm{w}} = \sum_c \frac{n_c}{N}\mathrm{F1}_c, \quad
\mathrm{F1}_{\mathrm{m}} = \frac{1}{|\mathcal{C}|}\sum_c \mathrm{F1}_c .
\end{equation}
We report both the weighted and the macro form throughout. On balanced test splits they coincide, and they do coincide to four decimal places in every row we report, which is the check that no result in this paper is produced by class skew. We also report the area under the ROC curve, because a threshold-free ranking measure can disagree with a thresholded one, and on DAIGT V2 it does.

\subsection{Paired statistics}
\label{sec:stats}

Every reported figure comes from a committed script that writes JSON, and every transformer baseline is read from the deployed checkpoint's own record rather than from the maximum over its hyperparameter grid. An earlier version of this analysis compared a grid maximum against a single fixed configuration and reported differences that did not exist, which is why baselines are checkpoint-matched throughout.

Seed spread over three seeds is reported as a range and only for the cell it was measured on, never borrowed across datasets or model families, since fine-tuning outcomes vary materially with initialisation and data order~\cite{Dodge2020fine,Mosbach2020stability}. A range over three runs is an order-of-magnitude indication and not a test statistic, so no comparison in this paper rests on one. We do not run parametric tests at $n \le 30$.

Where two models are compared on the same test set, the comparison is paired, because the two predictions are made on the same documents and treating them as independent samples discards that structure~\cite{Dietterich1998approximate}. Let $b$ be the number of documents the first model classifies correctly and the second does not, and $c$ the reverse. Under the null that the two have equal error, the discordant pairs split evenly, so
\begin{equation}
\label{eq:mcnemar}
b \mid (b+c)  \sim  \mathrm{Binomial}\!\left(b+c, \tfrac{1}{2}\right),
\end{equation}
and we report the exact two-sided binomial $p$-value~\cite{McNemar1947note} rather than the chi-squared approximation, which is unreliable when $b+c$ is small, as it is for several comparisons here.

An exact test reports whether a difference is distinguishable from zero, not how large it is, so each comparison also carries a paired bootstrap interval on the error difference~\cite{Efron1979bootstrap}. Writing $e^{(1)}_i$ and $e^{(2)}_i$ for the two models' per-document error indicators, we resample document indices with replacement $B=10{,}000$ times and take the empirical 2.5th and 97.5th percentiles of
\begin{equation}
\label{eq:boot}
\hat{\Delta}^{(b)}  =  \frac{1}{n}\sum_{i \in I_b}\Big(e^{(1)}_i - e^{(2)}_i\Big),
\end{equation}
where $I_b$ is the $b$th resample. Resampling index sets rather than each model's errors independently is what preserves the pairing.

A bootstrap resamples one test set and therefore says nothing about how the corpus was partitioned in the first place. Because the central claim of this paper is a null on HC3, and a null is exactly the kind of claim a single lucky partition can manufacture, the whole decomposition is repeated over five group-aware partitions of the same balanced sample, varying only the split seed. Section~\ref{sec:multisplit} reports what that shows.

\subsection{Controls}
\label{sec:controls}

Three controls guard the specific claims made below, and each was added after an uncontrolled version of the corresponding experiment produced a conclusion that did not survive.

The first is a \emph{pipeline-fires} control. Our whitespace cleaning produces bit-identical predictions for BERT on HC3, which taken alone is indistinguishable from a cleaning step that silently failed to run. We therefore apply the same pipeline to DAIGT V2, where it also removes emoji that WordPiece does represent, and confirm that predictions there do change. The pipeline fires, and the null is specific to whitespace.

The second is a \emph{tokenisation} control. Whether a detector can exploit a character-level cue is a property of its tokeniser and is directly testable, so we compare token identifier sequences for minimally different strings rather than reasoning about what the model probably does.

The third is a \emph{label-free} control. We score 400 inputs per condition that carry no human-versus-machine information at all, including uniform random character strings, token-shuffled text, punctuation-only strings, a repeated single token and the empty string. Without this control a confident prediction on corrupted input is indistinguishable from a learned response to it, which is the distinction Section~\ref{sec:labelfree} turns on.
